\section{Properties of Modules}

\subsection{Finitely generated, integral, and finite algebras}

\bdefn

    For an $A$-algebra $B$, and $S\subseteq B$ denote by $A[S]$ the smallest $A$-subalgebra containing $S$ in $B$.

\edefn

Let $\phi\colon A\to B$ be the algebra map, then $A[S]$ is the subring of $B$ generated by $\phi(A)\cup S$.
Alternatively, for every $s\in S$ define an indeterminant $x_s$, and let $X=\set{x_s}_{s\in S}$ be the set of all these indeterminants.
Consider the polynomial algebra $A[X]$, and define $f\colon A[X]\to B$ by $f(x_s)=s$.
Then $A[S]$ is just the image of $f$.
Explicitly, for a function $\alpha\colon S\to\bZ_{\geq0}$ with finite support, define $S^\alpha=\prod_{\alpha(s)\neq0}s^{\alpha(s)}$.
Then
$$ A[S] = \set{\hbox{Finite sums }\sum_\alpha\phi(a_\alpha)S^\alpha}[a_\alpha\in A] . $$

Note that we can consider $B$ as an $A$-module.
Then for $S\subseteq B$, we can consider the {\it submodule\/} generated by $S$: $\gen S$.
Since $A[S]$ is itself a module, $\gen S\subseteq A[S]$ automatically, and we do not have equality in general.

\bdefn

    An $A$-algebra $B$ is {\emphcolor finitely generated} if there is a finite $S\subseteq B$ such that $B=A[S]$.
    $B$ is a {\emphcolor finite algebra} if it is finitely generated as an $A$-module.

\edefn

A finite $A$-algebra is also finitely generated as an $A$-algebra, but not vice versa in general.

\bexam

    Suppose $S=\set s$ is a singleton, then $A[s]$ is the set of polynomials $\sum_k\phi(a_k)s^k$.

\eexam

\bexam

    Clearly the polynomial algebra $A[x_1,\dots,x_n]$ is a finitely generated $A$-algebra.
    But it is not a finite $A$-algebra: we can reduce this to the case of fields, in which $x_i^k$ are all linearly independent and thus
    there is no finite basis.

\eexam

\blemm

    An $A$-algebra $B$ is fg iff it is the quotient of some finite polynomial algebra $A[x_1,\dots,x_n]$ (up to isomorphism of $A$-algebras).

\elemm

\bproof

    This is equivalent to saying $B$ is fg iff there is a surjection $A[x_1,\dots,x_n]\to B$.
    This is clearly true: if $B$ is fg then for $s_1,\dots,s_n$ generating $B$ map $x_i\mapsto s_i$ and extend by freeness.
    Conversely, quotients of fg algebras are fg.
    \qed

\eproof

\bdefn

    Let $B$ be an $A$-algebra.
    An element $b\in B$ is {\emphcolor integral} over $A$ if $b$ is the root of a monic polynomial over $A$.
    That is, there exist $a_1,\dots,a_n\in A$ such that $b^n+a_1b^{n_1}+\cdots+a_n\cdot1=0$.

    $b\in B$ is {\emphcolor algebraic} over $A$ if $b$ is the root of any polynomial over $A$.

\edefn

Clearly integral elements are algebraic.
For fields the converse holds as well, as we can divide a polynomial by its leading coefficient.

\bexam

    The set of integral elements in the $\bZ$-algebra $\bQ$ is $\bZ$.
    Clearly an integer is integral, as $n\in\bZ$ is the root of the monic $x-n$.
    Conversely if $b=p/q\in\bQ$ is integral with $p,q$ coprime then
    $$ b^n + a_1b^{n-1} + \cdots + a_n1 = 0,\qquad a_i\in\bZ . $$
    Then we have that
    $$ p^n + a_1p^{n-1}q + \cdots + a_nq^n = 0 $$
    and so $q$ must divide $p^n$, but since $p,q$ are coprime this can only happen if $q=1$ and $b\in\bZ$.

\eexam

\blemm

    If $M$ is a fg $A$-module, then for every $\phi\in\endo_A(M)$ there exist $a_1,\dots,a_n\in A$ such that
    $$ \phi^n + a_1\phi^{n-1} + \cdots + a_n = 0 $$
    in $\endo_A(M)$.

\elemm

\bproof

    The proof is by a generalization of Cayley-Hamilton.
    Let $n\geq1$ and consider $R=\bZ[X_{ij}]$, the polynomial ring over $n^2$ indeterminants $X_{ij}$, $1\leq i,j\leq n$.
    Let $X=(X_{ij})\in M_n(R)$, and define its {\emphcolor characteristic polynomial}: $\Phi_X(T)\in R[T]$ given by
    $$ \Phi_X(T) = \det(T\id - X) . $$
    By Cayley-Hamilton over $R$'s field of fractions, $\Phi_X(X)=0$.

    Now if $A$ is any ring and $Y\in M_n(A)$, there is a unique ring morphism $\phi\colon R\to A$ which maps $X$ to $Y$ (coordinate-wise).
    Indeed, such a map must map $X_{ij}$ to $Y_{ij}$ which uniquely determines $\phi$.
    Then $\phi$ lifts to the ring morphism $R[T]\to A[T]$ (map $\sum_na_nT^n$ to $\sum_n\phi(a_n)T^n$), which maps $\Phi_X(T)$ to $\Phi_Y(T)$ (which is also given by the determinant $\det(T\id-Y)$ since the determinant
    is functorial: $\det(\phi(B))=\phi(\det(B))$).
    Thus $\Phi_X(X)=0$ implies that $\Phi_Y(Y)=\phi(\Phi_X(\phi(X)))=\phi(\Phi_X(X))=0$, as desired.
    So Cayley-Hamilton applies to any commutative ring.

    Now for our lemma, let $m_1,\dots,m_n\in M$ be generators.
    Then for all $j$ we have
    $$ \phi(m_j) = \sum_i a_{ij}m_i ,\qquad a_{ij}\in A . $$
    We associate to $\phi$ the representative matrix $[\phi]=(a_{ij})\in M_n(A)$ (which is not unique, unlike over fields).
    This has the usual properties that $[\phi^n]=[\phi]^n$ and $[\phi+\psi]=[\phi]+[\psi]$ (meaning that the composition or sum of representatives is a representative matrix).
    Now Cayley-Hamilton gives us that
    $$ \sum_n b_n[\phi]^n = 0 $$
    for some $b_n\in A$, and thus
    $$ \bracks{\sum_nb_n\phi^n} = 0 , $$
    and since only the zero endomorphism has the zero representative, this means $\sum_nb_n\phi^n=0$ as desired.
    \qed

\eproof

\bprop

    For an $A$-algebra $B$ and $b\in B$, the following are equivalent:
    \benum
        \item $b$ is integral over $A$,
        \item $A[b]\subseteq B$ is a finite algebra (fg module) over $A$,
        \item $A[b]$ is contained in a finite algebra over $A$,
        \item There exists a faithful $A[b]$-module $M$ such that $M$ is finitely generated as an $A$-module.
    \eenum

\eprop

\bproof

    Clearly $(2)$ implies $(3)$ implies $(4)$: take $M=C$, and since $A[b]\subseteq M$ it is faithful (the algebra map is inclusion, which is
    injective).
    $(1)$ implies $(2)$ since if $b^n+a_1b^{n-1}+\cdots+a_n1=0$ then $b^n$ is in $\gen{1,\dots,b^{n-1}}\subseteq B$ and then by induction so is
    every higher power of $b$, so that $A[b]=\gen{1,\dots,b^{n-1}}$ is finite.

    $(4)$ implies $(1)$: we have a ring morphism $\phi\colon A[b]\to\endo_A(M)$, so that every $f\in A[b]$ gives us an endomorphism
    $\phi(f)\in\endo_A(M)$.
    In particular by the above lemma there are $a_1,\dots,a_n\in A$ such that
    $$ \rho(b)^n + a_1\rho(b)^{n-1} + \cdots + a_n = 0 \in \endo_A(M) . $$
    For $f(b)=b^n+a_1b^{n-1}+\cdots+a_n\in A[b]$, we see that $\rho(f(b))=0$.
    But $M$ is a faithful $A[b]$-module, so $f(b)=0$ and $b$ is integral as desired.
    \qed

\eproof

\blemm

    \benum
        \item If $B$ is a finite $A$-algebra, and $M$ is a fg $B$-module, then $M$ is a fg $A$-module.
        \item If $B$ is a finite $A$-algebra, and $C$ is a finite $B$-algebra, then $C$ is a finite $A$-algebra.
    \eenum

\elemm

\bproof

    Clearly the second point follows from the first.
    Now if $M$ is generated by $m_1,\dots,m_n$ over $B$, and $B$ is generated by $b_1,\dots,b_k$ over $A$ then $b_im_j$ generates $M$ over $A$.
    Indeed, for $m\in M$ if $m=\sum_ib'_im_i$, then $b'_i=\sum_ja_{ij}b_j$ so
    $$ m = \sum_{i,j}a_{ij}(b_jm_i) , $$
    as desired.
    \qed

\eproof

\bcoro

    Let $B$ be an $A$-algebra, and suppose that $b_1,\dots,b_n\in B$ are integral over $A$.
    Then $A[b_1,\dots,b_n]$ is a finite $A$-algebra.

\ecoro

\bproof

    By induction: for $n=1$ this is our above proposition.
    Then if $b_{n+1}$ is integral over $A$, it is integral over $A[b_1,\dots,b_n]$ so that $A[b_1,\dots,b_{n+1}]$ is a finite $A[b_1,\dots,b_n]$-algebra, and thus a finite $A$-algebra
    by transitivity which we just showed.
    \qed

\eproof

\bcoro

    Let $B$ be an $A$-algebra, then
    $$ B_0 = \set{b\in B}[\hbox{$b$ is integral over $A$}] $$
    is an $A$-subalgebra of $B$

\ecoro

\bproof

    Clearly $0,1\in B_0$.
    Now if $b_1,b_2\in B_0$ then $A[b_1+b_2],A[b_1b_2]\subseteq A[b_1,b_2]$ which is a finite $A$-algebra and thus $b_1+b_2,b_1b_2$ are integral over $A$; $B_0$ is a subring of $B$.
    And if $a\in A,b\in B_0$ then $A[ab]\subseteq A[b]$ is finite, so that $ab\in B_0$.
    \qed

\eproof

\bdefn

    \benum
        \item For an $A$-algebra $B$, we define $A$'s {\emphcolor integral closure} in $B$ to be the set of all integral elements over $A$ in $B$; which we showed above is an $A$-subalgebra.
        \item An $A$-algebra $B$ is {\emphcolor integral} if every element of $B$ is integral over $A$.
        \item A morphism of rings $f\colon A\to B$ is {\emphcolor integral} if $B$ is an integral $A$-algebra via $f$.
        \item For a morphism of rings $f\colon A\to B$, $f(A)$ is {\emphcolor integrally closed} in $B$ if $B_0=f(A)$ (i.e.\ no $b\in B-f(A)$ is integral over $A$).
        \item A domain $A$ is {\emphcolor normal} if it is integrally closed in its field of fractions.
    \eenum

\edefn

\bexam

    We saw that $\bZ$ is normal (its field of fractions is $\bQ$); the same proof shows that any UFD is normal.
    In particular $k[x_1,\dots,x_n]$ and $\bZ[x_1,\dots,x_n]$ are normal.

    Consider the $k$-subalgebra $A=k[x^2,x^3]\subseteq k[x]$.
    Then $x\notin A$ is integral over $A$ since it is a root of $t^6-x^6\in A[t]$, and $x=x^3/x^2$ is in its field of fractions.
    Thus $A$ is not normal.

\eexam

\bnote

    There is similarity between integral morphisms and algebraic extensions of fields.
    But while algebraically closed fields exist, there are no ``integrally closed ring'': every ring can be embedded non-integrally into another ring.

\enote

\bcoro

    An $A$-alegbra $B$ is finite over $A$ iff it is integral and finitely generated.

\ecoro

\bproof

    If $B$ if a finite $A$-algebra, then for every $b\in B$, $A[b]\subseteq B$ is finite so $b$ is integral over $A$; $B$ is integral and finitely generated.
    Conversely if $B$ is integral and finitely generated, suppose $B=A[b_1,\dots,b_n]$; since each $b_i$ is integral we have that $B$ is finite.
    \qed

\eproof

\bcoro

    Let $B$ be an integral $A$-algebra, and $C$ be a $B$-algebra.
    \benum
        \item If $c\in C$ is integral over $B$, then it is integral over $A$.
        \item If $C$ is integral over $B$, then it is integral over $A$.
    \eenum

\ecoro

\bproof

    If $c$ is integral over $B$, then
    $$ c^n + b_1c^{n-1} + \cdots + b_n = 0 $$
    for some $b_i\in B$, which are integral over $A$.
    Then $B'=A[b_1,\dots,b_n]\subseteq B$ is a finite $A$-algebra, and $c$ is integral over $B'$ and so $B'[c]$ is a finite $B'$-algebra.
    Since $B'$ is a finite $A$-algebra, by transitivity $B'[c]$ is a finite $A$-algebra as well containing $A[c]$, so $c$ is integral over $A$.

    The second point follows from the first.
    \qed

\eproof

\bcoro

    Let $B$ be an $A$-algebra, and $B_0$ be the integral closure of $A$ inside $B$.
    Then $B_0$ is integrally closed in $B$.

\ecoro

\bproof

    If $b\in B$ is integral over $B_0$, then by the above corollary, since $B_0$ is an integral $A$-algebra, $b$ is integral over $A$ and thus $b\in B_0$.
    \qed

\eproof

Let $B$ be an $A$-algebra under $f\colon A\to B$.
\benum
    \item For an ideal $J\leq B$, define $J_A=f^{-1}(J)\leq A$.
    Then $f$ induces an injective homomorphism $A/J_A\to B/J$, so that $B/J$ is an $A/J_A$-algebra.
    \item For a multiplicative set $S\subseteq A$, $f(S)\subseteq B$ is multiplicative.
    Then we have a natural morphism $S^{-1}A\to f(S)^{-1}B$ mapping $a/s$ to $f(a)/f(s)$.
    Thus $f(S)^{-1}B$ is an $S^{-1}A$-algebra.
    \item Viewing $B$ as an $A$-module, we have already defined $S^{-1}B$, and this clearly coincides with $f(S)^{-1}B$.
    Thus $S^{-1}B$ is naturally an $S^{-1}A$-algebra.
\eenum

\bprop

    Let $B$ be an integral $A$-algebra.
    \benum
        \item Let $J\leq B$ be an ideal, then $B/J$ is an integral $A/J_A$-algebra.
        \item Let $S\subseteq B$ be multiplicative, then $S^{-1}B$ is an integral $S^{-1}A$-algebra.
    \eenum

\eprop

\bproof

    For $\bar x\in B/J$, lift it to $x\in B$, which is integral over $A$.
    Thus $x^n+a_1x^{n-1}+\cdots+a_n=0$ for $a_i\in A$, and then
    $$ \bar x^n + \bar a_1\bar x^{n-1} + \cdots + \bar a_n = 0 , $$
    where $\bar a_i=a_i+J_A\in A/J_A$, so that $\bar x$ is integral over $A/J_A$.

    An element of $S^{-1}B$ is of the form $\tilde x=x/f(s)$ with $x\in B,s\in S$.
    By assumption, $x^n+a_1x^{n-1}+\cdots+a_n=0$ and thus
    $$ \parens{\frac x{f(s)}}^n + \frac{a_1}s\parens{\frac x{f(s)}}^{n-1} + \cdots + \frac{a_n}{s^n} = \frac{x^n+a_1x^{n-1}+\cdots+a_0}{f(s)^n} = 0 , $$
    so that $\tilde x$ is integral over $S^{-1}A$ as desired.
    \qed

\eproof

\bprop

    Let $A\subseteq B$ be integral domains so that $B$ is integral over $A$.
    Then $A$ is a field iff $B$ is a field.

\eprop

\bproof

    If $A$ is a field, then let $0\neq b\in B$.
    It is integral over $A$ so $b^n+a_1b^{n-1}+\cdots+a_n=0$ for $a_i\in A$, and assume $n$ is minimal.
    Then
    $$ b(b^{n-1} + a_1b^{n-2} + \cdots + a_{n-1}) = -a_n . $$
    Since by minimality $b^{n-1}+a_1b^{n-1}+\cdots+a_{n-1}\neq0$, we have $a_n\neq0$.
    Thus $-a_n$ is a unit, and thus $b$ must be a unit as well (if $xy=u$ is a unit, then $x$ is a unit since $x(yu^{-1})=1$).
    So $B$ is a field.

    Conversely if $B$ is a field, then let $0\neq x\in A$.
    Then $x^{-1}\in B$ since it is a field.
    Since $B$ is integral over $A$, there are $a_i\in A$ such that
    $$ x^{-n} + a_1x^{-(n-1)} + \cdots + a_n = 0 , $$
    and multiplying by $x^{n-1}$ we get
    $$ x^{-1} + a_1 + \cdots + a_nx^{n-1} , $$
    so that $x^{-1}$ is the sum of elements from $A$, and thus is in $A$.
    \qed

\eproof

\bcoro

    Let $B$ be an integral $A$-algebra, and $\q\leq B$ a prime ideal.
    Let $\p=\q_A\leq A$ be its preimage, then $\q$ is maximal iff $\p$ is maximal.

\ecoro

\bproof

    Both $\p,\q$ are prime so that $B/\q$ and $A/\p$ are integral domains.
    As we showed, $B/\q$ is integral over $A/\p$.
    Thus $B/\q$ is a field iff $A/\p$ is a field; so $\q$ is maximal iff $\p$ is maximal.
    \qed

\eproof

\bcoro

    Let $B$ be an integral $A$-algebra, $\q\leq\q'\leq B$ be prime ideals such that $\q_A=\q'_A$.
    Then $\q=\q'$.

\ecoro

\bproof

    Let $\p=\q_A=\q'_A$; it is prime over $A$.
    Consider the following commutative square

    \commsq{A&B\cr A_\p&B_\p}
        {f}{f_\p}{\iota_A}{\iota_B}

    which commutes since the structure morphism $A\to S^{-1}A$ is a natural transformation (literally what the commutative square says).

    Since $f^{-1}\q=\p$ we see that $\q\cap f(A-\p)=\emptyset$ and similarly for $\q'$.
    Thus $\q B_\p\subseteq\q' B_\p\subseteq B_\p$ are prime ideals (since they don't intersect the localizing set $f(A-\p)$).

    We are going to show that $\q B_\p\subseteq B_\p$ is a maximal ideal so that $\q B_\p=\q' B_\p$.
    And then since $\q\mapsto\q B_\p$ is bijective with inverse $\iota_B^{-1}$, we will have that $\q=\q'$.

    Since $B$ is integral over $A$, $B_\p$ is integral over $A_\p$.
    Thus, by the above corollary, it is sufficient to show that $\tilde\p=f_\p^{-1}(\q B_\p)\subseteq A_\p$ is maximal.
    By assumption,
    $$ \iota_A^{-1}(\tilde\p) = \iota_A^{-1}(f_\p^{-1}(\q B_\p)) = f^{-1}\circ\iota_B^{-1}(\q B_\p) = f^{-1}(\q) = \p . $$
    Thus we have that $\tilde\p=\p A_\p$ (since $\iota_A^{-1}$ is bijective) which is maximal, as desired.
    \qed

\eproof

The main result we want is the following.

\bthrm

    Let $f\colon A\to B$ be an injective integral morphism.
    Then for every prime $\p\leq A$, there exists a prime ideal $\q\leq B$ such that $\p=\q_A$.

\ethrm

\bproof

    Since $B_\p=f(A-\p)^{-1}B$, and $0\notin f(A-\p)$ because $f$ is injective, $B_\p\neq0$.
    As before, $B_\p$ is integral over $A_\p$ and we have the same commutative square

    \commsq{A&B\cr A_\p&B_\p}
        {f}{f_\p}{\iota_A}{\iota_B}

    Let $\m\subseteq B_\p$ be maximal, so that $f_\p^{-1}(\m)\subseteq A_\p$ is maximal (as we showed $\q_A$ is maximal iff $\q$ is).
    Thus $f_\p^{-1}(\m)=\p A_\p$.
    Let $\q=\iota_B^{-1}(\m)$, which is prime in $B$, and
    $$ f^{-1}(\q) = f^{-1}\circ\iota_B^{-1}(\m) = \iota_A^{-1}\circ f_\p^{-1}(\m) = \iota_A^{-1}(\p A_\p) = \p , $$
    as desired.
    \qed

\eproof

\bnote

    Equivalently, this theorem says that if $f\colon A\to B$ is integral injective, then $f^*\colon\spec(B)\to\spec(A)$ is surjective.

    Note that injectivity is necessary here.
    This is because $\pi\colon A\to A/I$ is integral for any ideal $I\leq A$, but $\pi^*\colon\spec(A/I)\to\spec(A)$ is not surjective: its image is $V(I)$, the set of prime ideals containing $I$.

\enote

\bcoro

    Let $B$ be an integral $A$-algebra, then for every prime $\q\leq B$ and every prime ideal $\p'\supseteq\p=\q_A$, there exists a prime ideal $\q'\supseteq\q$ of $B$ such that $\q'_A=\p'$.

\ecoro

\bproof

    Consider the commutative square

    \commsq{A&B\cr A/\p&B/\q}
        {f}{\bar f}{\pi_A}{\pi_B} 

    Then $\p'\supseteq\p$ gives rise to a prime ideal $\bar\p'=\p'/\p\subseteq A/\p$.
    Since $B/\q$ is an integral $A/\p$-algebra, there is a prime ideal $\bar\q'\subseteq B/\q$ such that $\bar f^{-1}(\bar\q')=\bar\p'$.
    Then $\q'=\pi^{-1}_B(\bar\q')$ is a prime ideal of $B$ containing $\q$ and
    $$ f^{-1}(\q') = f^{-1}\circ\pi^{-1}_B(\bar\q') = \pi^{-1}_A\circ\bar f^{-1}(\bar\q') = \pi^{-1}_A(\bar\p') = \p' , $$
    as desired.
    \qed

\eproof

\bprop

    Let $B$ be an $A$-algebra, $C\subseteq B$ be the integral closure of $A$ inside of $B$.
    Then for any multiplicative $S\subseteq A$, $S^{-1}C$ is the integral closure of $S^{-1}A$ inside of $S^{-1}B$.

\eprop

\bproof

    Since $C$ is integral over $A$, $S^{-1}C$ is integral over $S^{-1}C$.
    Now if $b/s$ is integral over $S^{-1}A$, then $b/1=s\cdot b/s$ is integral over $S^{-1}A$ as well, and thus
    $$ \parens{\frac b1}^n + \frac{a_1}{s_1}\parens{\frac b1}^{n-1} + \cdots + \frac{a_n}{s_n} = 0 $$
    for $a_i/s_i\in S^{-1}A$.
    Let $t=s_1\cdots s_n$ and $u_i=t/s_i\in S$, and multiplying this equation by $t^n$ gives us
    $$ \parens{\frac{tb}1}^n + a_1u_1\parens{\frac{bt}1}^{n-1} + \cdots + a_nu_nt^{n-1} = 0 . $$
    Thus there is a $v\in S$ such that
    $$ v((bt)^n + a_1u_1(bt)^{n-1} + a_2u_2t(bt)^{n-2} + \cdots + a_nu_nt^{n-1}) = 0 . $$
    Then $btv$ is integral over $A$ (multiply by $v^{n-1}$), so that $btv\in C$.
    Then $b/s=(btv)/(stv)\in S^{-1}C$, as desired.
    \qed

\eproof

\bcoro

    Let $A$ be a domain, then the following are equivalent:
    \benum
        \item $A$ is normal,
        \item $S^{-1}A$ is normal for every multiplicative $S\subseteq A$ not containing $0$,
        \item $A_\p$ is normal for every prime $\p$,
        \item $A_\m$ is normal for every maximal $\m$.
    \eenum

\ecoro

\bproof

    Let $K$ be $A$'s field of fractions, and let $C\subseteq K$ be $A$'s integral closure in $K$.
    Then $S^{-1}C$ is $S^{-1}A$'s integral closure in $S^{-1}K=K$ for all multiplicative $S$ not containing $0$.
    If $A$ is normal, then $C=A$ so that $S^{-1}A$ is integrally closed in $K$, and thus normal.
    So (1) implies (2), which implies (3) which implies (4).

    Now if $A_\m$ is normal, then $A_\m\to C_\m$ is an isomorphism for all $\m$.
    We showed that this implies $A\to C$ is an isomorphism, so that $A$ is normal.
    \qed

\eproof

\subsection{Noetherian rings and modules}

\blemm

    Let $\Sigma$ be a poset, then the following conditions are equivalent:
    \benum
        \item every nonempty subset of $\Sigma$ has a maximal element,
        \item every increasing chain $x_1\leq x_2\leq\cdots$ in $\Sigma$ eventually stabilizes: at some point $x_n=x_{n+1}$ (equivalently, there are no infinite strictly
        increasing chains).
    \eenum
    A poset satisfying these equivalent conditions is said to have the {\emphcolor ascending chain condition} ({\emphcolor ACC}).

\elemm

\bproof

    The first implies the second as we can consider the subset $\set{x_i}$.
    The second implies the first since if $\Sigma'\subseteq\Sigma$ doesn't have a maximal element we can form a strictly increasing chain inductively.
    \qed

\eproof

\bdefn

    An $A$-module $M$ is {\emphcolor Noetherian} if its poset of $A$-submodules has the ACC.
    A ring $A$ is {\emphcolor Noetherian} if it is Noetherian as a module over itself (equivalently its ideals satisfy the ACC).

\edefn

\bexam

    Any field is a Noetherian ring, since it only has two ideals.

    $\bQ$ is not a Noetherian $\bZ$-module, since $\frac1{n!}\bZ$ does not stabilize (the reason for the factorial is so that it forms an increasing chain).

    The ring $k[x_1,x_2,\dots]$ with infinitely many indeterminants is not Noetherian, since $(x_1)\subsetneq(x_1,x_2)\subsetneq(x_1,x_2,x_3)\subsetneq\cdots$ forms a strictly increasing chain.

\eexam

\bprop

    An $A$-module is Noetherian iff every submodule is finitely generated.
    Thus a ring is Noetherian iff every ideal is finitely generated.

\eprop

\bproof

    If $M$ is Noetherian and $N\subseteq M$ is a submodule, suppose it is not finitely generated.
    Then inductively we will form an infinitely increasing sequence of submodules.
    Let $x_1\in N$ arbitrary, then since $N$ is not finitely generated $x_1,\dots,x_n$ does not generate $N$ so that there is an $x_{n+1}\in N-\gen{x_1,\dots,x_n}$.
    Then $N_n=\gen{x_1,\dots,x_n}$ is a strictly increasing infinite chain of submodules of $M$, a contradiction.

    Conversely if every submodule of $M$ is finitely generated, then $M$ is Noetherian.
    Otherwise, let $M_1\subsetneq M_2\subseteq\cdots$ be an increasing chain of submodules.
    Let $N=\bigcup_nM_n$ be their union, which is also a submodule of $M$, and thus is finitely generated.
    Suppose $x_1,\dots,x_n\in N$ generate $N$, and since $M_i$ form a chain there is an $m$ such that $x_1,\dots,x_n\in M_m$ and thus $M_m=N$ so that $M_i$ eventually stabilizes to $M_m$.
    \qed

\eproof

\bcoro

    Every PID is Noetherian.

\ecoro

\blemm

    Let $0\to M'\xto fM\xto gM''\to0$ be a short exact sequence of $A$-modules.
    Then $M$ is Noetherian iff $M',M''$ are.

    Equivalently, if $M$ is an $A$-module and $N\subseteq M$ a submodule, $M$ is Noetherian iff $N$ and $M/N$ are.

\elemm

\bproof

    If $M$ is Notherian, let $M'_n\subseteq M',M''_n\subseteq M''$ be increasing sequences of submodules.
    Then $f(M'_n)$ and $g^{-1}(M''_n)$ are submodules of $M$ and thus stabilize: $f(M'_n)=f(M'_{n+1})$ and $g^{-1}(M'_n)=g^{-1}(M'_{n+1})$ eventually.
    By the injectivity of $f$ and surjectivity of $g$, this means $M'_n$ and $M''_n$ both eventually stabilize.

    Conversely if $M',M''$ are Noetherian, let $M_n\subseteq M$ be an increasing sequence of submodules.
    Then $f^{-1}(M_n)=M_n\cap M'$ and $g(M_n)\subseteq M''_n$ are increasing sequences of submodules of $M',M''$ respectively.
    Thus eventually they stabilize: $f^{-1}(M_n)=f^{-1}(M_{n+1})$ and $g(M_n)=g(M_{n+1})$ eventually.
    We claim that this means $M_n=M_{n+1}$.

    Indeed, for $x\in M_{n+1}$ consider $g(x)\in g(M_{n+1})$ so there is a $y\in M_n$ such that $g(x)=g(y)$.
    Then $x-y\in\ker g=\im f$ so that $x-y=f(z)$ for $z\in M'$, and $x-y\in M_{n+1}$ so that $z\in f^{-1}(M_{n+1})$ and thus $z\in f^{-1}(M_n)$.
    This means that $x-y\in M_n$ and thus $x\in M_n$ as desired.
    \qed

\eproof

\bcoro

    \benum
        \item The finite direct sums of Noetherian modules is Noetherian.
        \item If $A$ is a Noetherian ring and $M$ is a fg $A$-module, then $M$ is a Notherian module.
        \item If $A$ is a Notherian ring and $B$ a finite $A$-algebra, then $B$ is a Noetherian ring.
        In particular, $A/I$ is Notherian for any ideal $I\leq A$.
    \eenum

\ecoro

\bproof

    For the first point it is sufficient to show this for binary direct sums.
    If $M,N$ are Noetherian, then $0\to M\to M\oplus N\to N\to0$ is a short exact sequence and thus $M\oplus N$ is Notherian by above.

    Since $M$ is fg, $M\cong A^n/N$ for some submodule $N\leq A^n$ (consider the surjection $A^n\to M$ mapping to generators).
    By the first point, $A^n$ is Noetherian and thus $M$ is as well.

    $B$ is a fg $A$-module, and thus a Noetherian $A$-module.
    Since ideals of $B$ are $A$-submodules, the ACC on $A$-submodules implies ACC on $B$'s ideals.
    \qed

\eproof

\blemm

    Let $A$ be a Noetherian ring, then all of its localizations are Noetherian rings as well.

\elemm

\bproof

    Let $S\subseteq A$ be multiplicative, and consider an increasing sequence of ideals $\tilde I_1\subseteq\tilde I_2\subseteq\cdots$ in $S^{-1}A$.
    Then $I_n=\iota^{-1}(\tilde I_n)$ forms an increasing sequence of ideals of $A$, and thus stabilizes.
    Since $\tilde I_n=S^{-1}I_n$, we get that $\tilde I_n$ stabilizes.
    \qed

\eproof

\bthrm[title=Hilbert basis theorem]

    If $A$ is Noetherian, then $A[x]$ is a Noetherian ring.

\ethrm

\bproof

    Let $I\subseteq A[x]$ be an ideal, we wish to show it is finitely generated.
    Let us adopt the notation that for $f\in A[x]$, $a_f\in A$ denotes the leading term of $f$ (and $a_f=0$ only when $f=0$).

    We claim that $J=\set{a_f}[f\in I]$ is an ideal in $A$.
    If $f\in I$ and $b\in A$ then either $ba_f=0$ or $ba_f=a_{bf}$; in both cases $ba_f\in J$.
    Now if $f,g\neq0$ are in $I$ then $a_f+a_g$ is either zero or the leading term of $x^{\deg g}f+x^{\deg f}g\in J$ and thus $a_f+a_g\in J$.

    Since $J$ is Noetherian, there are $f_1,\dots,f_n\in I$ such that $J=(a_{f_1},\dots,a_{f_n})$.
    Let $d_i=\deg f_i$, and $d$ be the maximum $d_i$.
    Consider the fg $A$-submodule
    $$ M = \sum_{i=0}^{d-1}Ax^i \subseteq A[x] . $$
    Then $M\cap I$ is an $A$-submodule of $M$, which is Noetherian (since $M$ is fg and $A$ is Noetherian so $M$ is Noetherian, and thus its submodules are Noetherian).

    Now define $I'=(f_1,\dots,f_n)\subseteq I$.
    We claim that
    $$ I = I' + (M\cap I)\hbox{ as $A$-submodules} . $$
    Let $f\in I$, then we claim there is an $f'\in I'$ such that $\deg(f-f')<d$ and thus $f-f'\in M\cap I$.
    If $\deg f<d$ then we are done, so assume $\deg f\geq d$.
    By induction, it is sufficient to show that there is an $f'\in I'$ such that $\deg(f-f')<\deg f$.

    Now, $a_f\in J=(a_{f_1},\dots,a_{f_n})$ so that $a_f=\sum_ib_ia_{f_i}$ for $b_i\in A$.
    Then define
    $$ f' = \sum_{i=1}^n b_ix^{(\deg f)-d_i}f_i \in I' . $$
    By assumption, $\deg f\geq d\geq d_i$ so that $\deg(x^{(\deg f)-d_i}f_i)=d_i$, and so $\deg f'\leq\deg f$.
    Moreover, the leading coefficient of $f'$ is
    $$ a_{f'} = \sum_{i=1}^n b_ia_{f_i} = a_f, $$
    so that $\deg(f-f')<\deg f$ as desired.

    Now since $I=I'+(M\cap I)$, $I'$ is a fg ideal and $M\cap I$ is a fg $A$-module (and thus fg by $A[x]$), we have that $I$ is a fg ideal, as desired.
    \qed

\eproof

\bcoro

    \benum
        \item If $A$ is a Noetherian ring, so is $A[x_1,\dots,x_n]$.
        \item If $A$ is a Noetherian ring, and $B$ is a fg $A$-algebra, $B$ is Noetherian.
        \item Any fg algebra over a field is Noetherian.
        \item Any fg algebra over a PID (in particular $\bZ$) is Noetherian.
    \eenum

\ecoro

\bproof

    The first point follows by induction and Hilbert's basis theorem.
    The second point is because if $B$ is a fg $A$-algebra, it is isomorphic to $A[x_1,\dots,x_n]/I$ for some ideal $I\leq A[x_1,\dots,x_n]$, and is thus Noetherian.
    The last two points are due to the second (fields and PIDs are Noetherian).
    \qed

\eproof

We now utilize these results to reprove a lemma we showed above.

\blemm

    Let $A$ be a ring and $T\in M_n(A)$ a matrix.
    Then there are $a_1,\dots,a_n\in A$ such that
    $$ T^n + a_1T^{n-1} + \cdots + a_m = 0 . $$

\elemm

\bproof

    Denote $T=(a_{ij})$, and consider the subalgebra $A'=\bZ[\set{a_{ij}}]\subseteq A$.
    This is a fg $\bZ$-algebra and thus Noetherian.
    Since $T\in M_n(A')$, we can assume from the getgo that $A$ is Noetherian.

    Since $M_n(A)\cong A^{n^2}$ as $A$-modules, $M_n(A)$ is a Noetherian $A$-module.
    Thus the increasing sequence of $A$-modules
    $$ N_m = \gen{1,T,\dots,T^m} \subseteq M_n(A) $$
    stabilizes, so eventually $N_n=N_{n-1}$ so that $T^n\in N_{n-1}$ and we get the desired result.
    \qed

\eproof

\blemm

    If $M$ is a fg $A$-module, then for $\phi\in\endo_A(M)$, there are $a_1,\dots,a_n\in A$ such that
    $$ \phi^n + a_1\phi^{n-1} + \cdots + a_n = 0 . $$

\elemm

\bproof

    When $M=A^n$ there is a natural isomorphism of $A$-algebras $M_n(A)\cong\endo_A(M)$, so the result follow from the above lemma.
    Otherwise, consider a surjective map $\pi\colon A^n\to M$, then we claim there is a $\tilde\phi\colon A^n\to A^n$ such that the below diagram commutes

    \commsq{A^n&A^n\cr M&M}
        {\tilde\phi}{\phi}{\pi}{\pi}

    Meaning $\pi\circ\tilde\phi=\phi\circ\pi$.
    For every $i$ choose a preimage $V_i\in\pi^{-1}(\phi(\pi(e_i)))\in A^n$ and define
    $$ \tilde\phi(x_1,\dots,x_n) = \sum_{i=1}^n x_iv_i . $$
    This satisfies the condition.

    Now there is a monic polynomial $f\in A[x]$ such that $f(\tilde\phi)=0$ by the $M=A^n$ case.
    Then
    $$ f(\phi)\circ\pi = \pi\circ f(\tilde\phi) = 0, $$
    which implies $f(\phi)=0$ because $\pi$ is surjective.
    \qed

\eproof

\subsection{Simple modules and length}

\bdefn

    A poset $\Sigma$ has the {\emphcolor descending chain condition} ({\emphcolor DCC}) if it satisfies the dual of the ACC (equivalently, $\Sigma^\op$ has ACC).
    This means that every subset has a minimal element, or every descending chain stabilizes.

\edefn

\bdefn

    An $A$-module is {\emphcolor Artinian} if it has DCC on its submodules.
    A ring $A$ is {\emphcolor Artinian} if it is Artinian as a module over itself, equivalently it has DCC on its ideals.

\edefn

Note that the short exact sequence result on Noetherian rings holds for Artinian rings as well.
Thus finite direct sums of Artinian modules are Artinian, fg modules over Artinian rings are Artinian, and thus finite algebras over Artinian rings are Artinian.

Note that it is not the case that a module is Artinian iff its submodules are all fg.
That is, Artinian is not equivalent to Noetherian (and for modules they are distinct conditions, for rings we will see that Artinian implies Noetherian).

\bexam

    $\bZ$ is a Noetherian but not an Artinian ring.
    It is not Artinian because $\bZ\supsetneq(n)\supsetneq(n^2)\supsetneq(n^3)\supsetneq\cdots$ is an infinite decreasing sequence of ideals.

    Similarly $k[x]$ is Noetherian but not Artinian: $(x^k)$ forms a descending sequence of ideals.

\eexam

\bdefn

    A nonzero module $M$ is {\emphcolor simple} if it has no non-trivial submodules.

\edefn

\blemm

    An $A$-module $M$ is simple iff it is isomorphic (as $A$-modules) to $A/\m$ for a maximal ideal $\m\leq A$.

\elemm

\bproof

    If $M$ is simple, then for $0\neq m\in M$ the submodule $\gen m\subseteq M$ is isomorphic to $A/\ann(m)$ (the map $a\mapsto am$ has kernel $\ann(m)$).
    Since $M$ is simple, $\gen m=M$ and so $A/\ann(m)\cong M$.
    There is then a bijection between submodules of $M$ and ideals of $A$ containing $\ann(m)$.
    Since $M$ is simple, there cannot be any nontrivial such ideals, so $\ann(m)$ is maximal.

    And if $M\cong A/\m$ then again there is a bijection between submodules of $M$ and ideals containing $\m$.
    So $M$ is simple, as $\m$ is maximal.
    \qed

\eproof

\bdefn

    A {\emphcolor composition series} of $M$ is a chain of submodules
    $$ 0 = M_0 \subsetneq M_1 \subsetneq M_2 \subsetneq \cdots \subsetneq M_n = M $$
    which is maximal.
    This is equivalent to $M_{i+1}/M_i$ being simple for all $i$ (no modules between $M_i$ and $M_{i+1}$).
    The {\emphcolor length} of such a series is $n$ (not $n+1$).

    Denote by $\ell_A(M)$ the minimal length of a composition series of $M$ (or infinity if no composition series exists).

\edefn

\blemm

    If $M$ has finite length $\ell(M)=n$, then all of its composition series have the same length $n$, and every chain of submodules can be extended to a composition series.

\elemm

\bproof

    First we claim that if $N\subsetneq M$ is a submodule, then $\ell(N)<\ell(M)$.
    Indeed, if $\set{M_i}$ is a composition series for $N$, then $N_i=M_i\cap N$ is a sequence of submodules of $N$.
    Then $N_i=N_{i+1}\cap M_i$ so
    $$ N_{i+1}/N_i = N_{i+1}/(N_{i+1}\cap M_i) \cong (N_{i+1}+M_i)/M_i \subseteq M_{i+1}/M_i $$
    which is either simple or zero.
    Thus removing repeated values, we have formed a composition series for $N$ of shorter length, so $\ell(N)\leq\ell(M)$.

    Now if $\ell(N)=\ell(M)$ then $N_{i+1}+M_i=M_{i+1}$ for all $i$, and thus we claim $N_i=M_i$ for all $i$ so that $N=M$.
    This is by induction, since $M_0=N_0=0$, and $N_{i+1}=N_{i+1}+N_i=N_{i+1}+M_i=M_{i+1}$, as desired.

    Now if $\set{M_i}_0^n$ is any chain in $M$, then
    $$ \ell(M) = \ell(M_n) > \ell(M_{n-1}) > \cdots > \ell(M_0) = 0, $$
    so $\ell(M)\geq n$.

    Now if $\set{M_i}$ is any composition series of $M$ of length $n$, then $\ell(M)\geq n$ by what we just showed.
    But by definition $\ell(M)\leq n$ so that $\ell(M)=n$ and thus all composition series have length $\ell(M)$.

    And for any general chain of submodules of $M$, if it is not a composition series then new terms can be inserted.
    This process can be repeated, and must terminate (since once it reaches a length of $\ell(M)$ it must be a composition series).
    \qed

\eproof

\bnote

    Note that we can prove something stronger, the {\emphcolor Jordan-Holder theorem}: for any two composition series $M_i,N_j$ there is a permutation $\sigma$ between $i$ and $j$
    such that $M_{i+1}/M_i\cong N_{\sigma(i)+1}/N_{\sigma(i)}$.
    A proof of this can be found in my representation theory homework somewhere.

\enote

\blemm

    An $A$-module $M$ has finite length iff it is both Noetherian and Artinian.

\elemm

\bproof

    If $M$ has finite length, then any increasing/decreasing sequence has length bound by $\ell(M)$ so that $M$ is both Noetherian and Artinian.

    Conversley if $M$ is Artinian and Noetherian, then consider the set of submodules $N\subseteq M$ of finite length.
    It is nonempty since it contains $0$, and because $M$ is Noetherian contains a maximal element $M'$.
    If $M=M'$ then we are done.
    Otherwise, the set of submodules $M'\subsetneq N\subseteq M$ is nonempty (containing $M$) and has a minimal element $N$ since $M$ is Artinian.
    But then $N/M'$ is simple, and so composing a composition series of $M'$ with $\subsetneq N$ forms a composition series of $N$, so $N$ has finite length, contradicting $M'$'s maximality.
    \qed

\eproof

\bprop

    Let $M$ be fg over a Notherian ring $A$.
    \benum
        \item If $M\neq0$, then $M$ has a submodules $N\subseteq M$ such that $N\cong A/\p$ for some prime $\p\leq A$.
        \item There exists a chain $0=M_0\subsetneq M_1\subsetneq\cdots\subsetneq M_n=M$ of submodules such that for each $i$, there is an isomorphism of $A$-modules
        $M_{i+1}/M_i\cong A/\p_i$ for prime ideals $\p_i\leq A$.
        \item The chain in the previous point is not necessarily unique, but
        \benum
            \item for a prime ideal $\p\leq A$ we have $\ann M\subseteq\p$ iff $\p_i\subseteq\p$ for some $i$.
            In particular, $\ann M\subseteq\p_i$ for all $i$, and the minimal ideals between $\p_1,\dots,\p_n$ are precisely minimal prime ideals $\ann(M)\subseteq\p$.
            \item for every minimal prime ideal $\ann(M)\subseteq\p$, the $\A_\p$ module $M_\p$ has finite length equal to the number of times $\p$ occurs in $\set{\p_1,\dots,\p_n}$.
        \eenum
    \eenum

\eprop

\bproof

    Let $\Theta$ be the set of all ideal $I\leq A$ of the form $I=\ann(m)$ for $0\neq m\in M$.
    Since $A$ is Noetherian, there is an $m$ such that $\ann(m)$ is maximal.
    Since $\gen m\subseteq M$ is isomorphic to $A/\ann(m)$, it is sufficient to show $\ann(m)$ is prime.
    So suppose $a,b\notin\ann(m)$, then $bm\neq0$ and $\ann(m)\subseteq\ann(bm)$.
    Thus by maximality $\ann(m)=\ann(bm)$ and then $a\notin\ann(bm)$ so $abm\neq0$ meaning $ab\notin\ann(m)$, showing it is prime.
    This proves the first point.

    Let $\Sigma$ be the set of all submodules of $M$ which possess such a chain.
    It is nonempty since the zero module is in it.
    Since $M$ is Noetherian, let $M'$ be a maximal element of $\Sigma$.

    Now consider $M''=M/M'$; if $M''=0$ we are finished.
    Otherwise by the first point it has a submodule $N''=N'/M'$ isomorphic to $A/\p$.
    But then composing the chain of $M'$ with $\subsetneq N'$ gives such a chain, contradicting $M'$'s maximality.

    For the final point, notice that for $M'\subsetneq M$,
    $$ \ann(M')\cdot\ann(M/M') \subseteq \ann(M) \subseteq \ann(M')\cap\ann(M/M') . $$
    The second inclusion is clear, the first is because for $a\in\ann(M')$ and $b\in\ann(M/M')$ and $m\in M$ then $bm\in M'$ and $abm\in aM'=0$.
    Thus for a prime ideal $\p\leq A$, $\ann(M)\subseteq\p$ iff $\ann(M')\subseteq\p$ or $\ann(M/M')\subseteq\p$.
    Thus by induction $\ann(M)\subseteq\p$ iff $\ann(M_{i+1}/M_i)\subseteq\p$ for some $i$.
    Since $\ann(A/\p_i)=\p_i$, the first point follows.

    For the second point, the chain $\set{M_i}$ gives the chain $\set{S^{-1}M_i}$ and
    $$ S^{-1}M_{i+1}/S^{-1}M_i \cong S^{-1}(M_{i+1}/M_i) \cong S^{-1}(A/\p_i) \cong S^{-1}A/S^{-1}\p_i . $$
    Moreover $S^{-1}\p_i=S^{-1}A$ iff $S\cap\p_i\neq\emptyset$, and otherwise is a prime ideal.
    For $S=A-\p$ for a minimal prime ideal $\ann(M)\subseteq\p$, then $S^{-1}\p_i=A_\p$ if $\p\neq\p_i$ and otherwise $S^{-1}\p_i=\p A_\p$ is a maximal ideal.
    Thus the sequence $(M_i)_\p$ forms a composition series after deleting the repeated entries (when $\p\neq\p_i$), and we are left with a composition series (since the quotients are
    quotients by a maximal ideal) whose length is equal to the number of $\p_i$ equal to $\p$.
    \qed

\eproof

\bdefn

    A minimal prime ideal is a prime ideal not containing any other prime ideal.

\edefn

\blemm

    A Noetherian ring $A$ has finitely many minimal prime ideals.

\elemm

\bproof

    The minimal prime ideals of $A$ are the minimal primes containing $\ann(A)=0$.
    Then for any chain as in the above lemma (whose components are $A/\p$ for some prime), the minimal prime ideals are the minimal primes from $\set{\p_1,\dots,\p_n}$ which is finite.
    \qed

\eproof

\bprop

    A ring is Artinian iff it is Noetherian and every prime ideal is maximal.

\eprop

\bproof

    If $A$ is Noetherian and every prime ideal is maximal, then every prime ideal is also minimal.
    Thus $A$ has finitely many prime/maximal ideals, $\m_1,\dots,\m_k$.
    Then their product $I=\m_1\cdots\m_k$ is contained in the intersection of all prime ideals, $\nil(A)$.
    Because $A$ is Noetherian $I$ is finitely generated, so $I^n=0$ for some $n$ (since every element of $I$ is nilpotent, and it is finitely generated).

    We have a chain of submodules
    $$ 0 = M_0 \subseteq M_1 \subseteq \cdots \subseteq M_k = A , $$
    where $M_{i+1}/M_i\cong A/\m_i$ for some $i$ (by the above proposition).
    As a subquotient, $M_{i+1}/M_i$ is a Noetherian $A$-module, and thus a Notherian module over the field $K=A/\m_i$.
    A Noetherian vector space is finite-dimensional (it is fg), so that $M_{i+1}/M_i$ is of finite length over $K$ (its length is equal to its dimension).
    Then $M_{i+1}/M_i$ has finite length over $A$ as well, and then inductively we see that each $M_i$ has finite length, and in particular $A$ is Artinian.

    Conversely if $A$ is Artinian, then all of its prime ideals are maximal (homework), so we need only show it is Noetherian.
    If $\m_1,\dots,\m_{n+1}$ are distinct maximal ideals then $\m_1\cdots\m_n\supset\m_1\cdots\m_{n+1}$ is strict, so by Artinian-ness there must be finitely many maximal ideals.
    If the inclusion were not strict then $\m_1\cdots\m_n\subseteq\m_{n+1}$, and then by $\m_{n+1}$'s primality, $\m_i\subseteq\m_{n+1}$ for some $i$.
    But then $\m_i=\m_{n+1}$ by maximality, a contradiction to distinctness.

    So let $I=\m_1\cdots\m_n$ be the product of all $A$'s prime/maximal ideals.
    Consider the descending chain $I\supseteq I^2\supseteq I^3\supseteq\cdots$ which must terminate, so $I^k=I^{k+1}$ eventually.
    We claim that $I^k=0$.

    If not, let $x_1,\dots,x_{k+1}\in I$ such that $x_1\cdots x_{k+1}\neq 0$.
    Because $I^k=I^{k+1}$, $\prod_{i=2}^{k+1}x_i=\sum_\alpha y_{\alpha,1}\cdots y_{\alpha,k+1}$ where the sum is over a finite set and each $y_{\alpha,i}\in I$.
    Thus $x_1y_{\alpha_0,1}\cdots y_{\alpha_0,k+1}\neq0$ for some $\alpha_0$.
    Then from the original tuple $x_1,\dots,x_{k+1}\in I$ we obtain a new tuple $x'_1=x_1,x'_2,\dots,x'_{k+2}$ ($x'_i=y_{\alpha_0,i-1}$) with nonzero product, and the first element unchanged.

    We similarly may write the product of the final $k$ elements $x'_3\cdots x'_{k+2}$ as an element of $I^{k+1}$ and obtain $x_1,x'_2,x'_3,x''_4,\dots,x''_{k+3}$ with nonzero product.
    Continuing this inductively, we produce an infinite sequence $y_1,y_2,\dots$ such that $y_1\cdots y_m\neq0$ for all $m>0$.
    Now consider the descending chain $(y_1)\supseteq(y_1y_2)\supseteq\cdots$, which stabilizes since $A$ is Artinian.
    Thus for some $m>0$ and $a\in A$,
    $$ \prod_{i=1}^my_i = a\prod_{i=1}^{m+1}y_i \implies \parens{\prod_{i=1}^my_i}(1-ay_{m+1}) = 0 . $$
    Since $y_{m+1}\in I$ and thus $ay_{m+1}\in I$ is in the Jacobson radical ($I$ is the product of all maximal ideals, and thus contained in their intersection), we have that $1-ay_{m+1}\in A^\times$.
    Thus $\prod_{i=1}^my_i=0$, a contradiction.

    So $I^k=0$, and so there is a chain (taking $M_j=\m_1\cdots\m_{n-j}$),
    $$ 0 = M_0 \subseteq M_1 \subseteq \cdots \subseteq M_m = A $$
    where $M_{i+1}/M_i$ is annihlated by some $\m_i$, and thus is Artinian over the field $K=A/\m_i$ so finite-dimensional over $A$.
    Thus $M_{i+1}/M_i$ has finite length over $A/\m_i$, and thus over $A$, and then by induction we see that each $M_i$ has finite length, in particular so does $A$ and thus is Noetherian.
    \qed

\eproof

\bnote

    Recall that $\bZ$ is Notherian but not Artinian.
    This is because while all of its non-zero prime ideals are maximal, $(0)$ is a prime ideal which is not maximal.
    We see that in general an integral domain which is not a field cannot be Artinian.

\enote


